
\section{Methods}

Maximum of 1200 words. 

The method describes the entire approach, so that others can in principle reproduce the work. The description of the method takes place at 3 levels. 

First, the chosen overall method is described, in line with the objective of the assignment. Within the HBO-ICT program for example, assignments often consist of creating a (partly) innovative artifact (product or service). Design Science Research (DSR) is the standard method within the program for such design-oriented research. 

Subsequently, the concrete steps carried out within the chosen method are described in the form of phases and/or activities. It becomes clear whether these were carried out sequentially, parallel and/or iteratively. An illustration is highly recommended here. Specific from the DSR method should in any case be included: identifying requirements, carrying out grounding, designing, developing and testing the product or service and
disseminating the results of the work carried out. The field test from the DSR framework is optional, as a rule students will not arrive at a product that is ripe enough for a field test within the boundaries of the education.

Last but not least, the methods/methods used specifically within the phase(s) / activity(ies) are described concretely and in depth, where possible consisting of best or good practices. It is always important that the choices made are substantiated in relation to the context and objective of the assignment
